Soit une fonction $f$ continue sur $]a,b[$ ($-\infty \leq a < b \leq +\infty$). On s'int�resse ici � la meilleure approximation polynomiale de $f$ selon une norme discr�te. Pour cela on pr�l�ve $m$ couples de points sur le graphe de $f$ :

\[
(x_i, y_i)_{1 \leq i \leq m},
\]

avec $x_i \in ]a,b[$ et $y_i = f(x_i)$. On va chercher le polyn\^ome $p_n$ de degr� $n$ qui approche au mieux $f$ en ces $m$ points, autrement dit qui soit le minimum de l'erreur :

\[
E(q_n) = \sum_{i=1}^m | y_i - q_n(x_i) |^2,
\]

pour tout $q_n \in \mathcal{P}_n$. On supposera que le degr� $n$ du polyn\^ome d'approximation est bien inf�rieur au nombre $m$ de points $(x_i,y_i)$.\\

On va programmer ici la m�thode de d�termination du meilleur polyn�me � l'aide du logiciel SAGE
({\tt www.sagemath.org}). On va consid�rer la fonction :

\[
f : x \rightarrow \sin ( 2\pi \cos(\pi x) ),
\]

restreinte � l'intervalle $[0,1]$. On choisira les points $x_i$ uniform�ment.

\begin{enumerate}

\item D�finir la fonction $f$ ci-dessus sous SAGE, et tracer $f$ sur l'intervalle $[0,1]$.

\item Pour $m$ donn�, stocker dans des listes les $m$ points $(x_i,y_i)$ n�cessaires � l'approximation polynomiale.

\item Montrer que minimiser $E(q_n)$ sur l'ensemble des polyn�mes de degr� $n$ revient �
r�soudre un syst�me matriciel de la forme :

\[
\mathbf{M}^T\mathbf{M}\:\mathbf{a} = \mathbf{M}^T \mathbf{y},
\]

avec $\mathbf{a}$ vecteur colonne qui contient tous les coefficients du polyn�me de meilleure approximation $p_n$.

\item Construire sous SAGE la matrice $\mathbf{M}$.

\item R�soudre le syst�me lin�aire ci-dessus pour obtenir les coefficients $(a_j)$ de $p_n$.

\item Tracer le polyn�me $p_n$ superpos� � $f$.

\item Calculer l'erreur (discr�te) quadratique entre $p_n$ et $f$.

\item Augmenter le degr� $n$ du polyn\^ome et d\'eterminer la valeur de $n$ \`a partir de laquelle l'erreur vous semble acceptable.

\end{enumerate}